\chapter{Ölçek Puanlarının Güvenirliği}
\label{ch:reliability}
\index{güvenirlik}
\index{ölçek!güvenirlik}

\begin{hedefler}
Bu bölümün sonunda okuyucu:
\begin{itemize}
  \item güvenirliği bir ölçeğin değişmez özelliği değil, belirli bir örneklemde
  ve uygulama koşulunda elde edilen puanlara ilişkin kanıt olarak açıklayabilecek,
  \item güvenirlik ile geçerliği birbirinden ayırabilecek,
  \item iç tutarlılık, test--tekrar test ve değerlendiriciler arası güvenirlik
  sorularını uygun katsayılarla eşleştirebilecek,
  \item ters maddeleri ve eksik veriyi analizden önce denetleyebilecek,
  \item Cronbach alfa katsayısını hesaplayıp varsayımları ve sınırlılıklarıyla
  yorumlayabilecek,
  \item düzeltilmiş madde--toplam korelasyonu ile ``madde silinirse alfa''
  çıktısını mekanik karar vermeden değerlendirebilecek,
  \item omega katsayısının alfa karşısındaki temel gerekçesini açıklayabilecek,
  \item Python, R ve SPSS çıktılarından ölçülü bir güvenirlik paragrafı
  yazabilecektir.
\end{itemize}
\end{hedefler}

\begin{hazirlik}
Beş maddelik bir ölçeğin alfa katsayısı 0,92 bulunmuştur. Bu sayı ölçeğin tek
boyutlu, geçerli ve her grupta güvenilir olduğunu tek başına gösterir mi?
Gerekçenizi üç ayrı cümleyle yazın.
\end{hazirlik}

\section{PDR vakası: Yardım arama tutumu ölçeği}

Bir psikolojik danışma merkezi, öğrencilerin profesyonel yardım aramaya yönelik
tutumlarını beş maddelik bir ölçekle değerlendirmektedir. Maddeler 1
(``Kesinlikle katılmıyorum'') ile 5 (``Kesinlikle katılıyorum'') arasında
puanlanmıştır. Dördüncü madde olumsuz yöndedir:

\begin{quote}
``Sorunlarımı bir uzmana anlatmanın yararlı olacağını düşünmem.''
\end{quote}

Bu madde ters puanlanmadan toplam puana eklenirse yüksek yardım arama tutumu
gösteren öğrenciler dördüncü maddede düşük puan alacaktır. Böylece aynı yapıyı
ölçmesi amaçlanan maddeler yapay biçimde ters yönde ilişkilendirilebilir.
Güvenirlik analizi bu nedenle bir menü komutuyla değil; madde yönü, puanlama
anahtarı, hedef grup ve ölçüm amacı incelenerek başlar.

\begin{researchbox}
``Ölçeğin güvenirliği 0,80'dir'' cümlesi eksiktir. Katsayı; hangi puanın,
hangi örneklemde, hangi uygulama koşulunda ve hangi yöntemle değerlendirildiği
belirtilerek raporlanmalıdır.
\end{researchbox}

\section{Gözlenen puan, hedef puan ve ölçme hatası}
\index{klasik test kuramı}
\index{ölçme hatası}

Klasik test kuramının temel gösterimi
\[
X=T+E
\]
biçimindedir. Burada $X$ gözlenen puan, $T$ belirli ölçme koşullarında hedeflenen
puan bileşeni ve $E$ ölçme hatasıdır. Aynı bireyin gözlenen puanı; madde
örneklemi, uygulama ortamı, zaman, değerlendirici veya geçici durumlar nedeniyle
değişebilir.

\begin{definitionbox}
\textbf{Güvenirlik}, belirli bir puan yorumunda gözlenen farklılıkların ne
ölçüde hedeflenen puan farklılıklarını yansıttığına ilişkin kanıttır. Katsayı,
ölçeğin değişmez etiketi değil; puanların belirli kullanımına ilişkin bir
özettir.
\end{definitionbox}

Basitleştirilmiş klasik gösterimde güvenirlik
\[
\rho_{XX'}=\frac{\Var(T)}{\Var(X)}
=1-\frac{\Var(E)}{\Var(X)}
\]
olarak yazılır. Hata değişkenliği büyüdükçe aynı gözlenen puan farkının hedef
puan farkını temsil etme derecesi azalır. Bu formül her hata kaynağını tek tek
tanımlamaz; hangi kaynakların hata sayılacağı ölçüm tasarımına bağlıdır.

\section{Güvenirlik ve geçerlik aynı soru değildir}
\index{geçerlik}

Güvenirlik puanların tutarlılığı ve ölçme hatasıyla, geçerlik ise puanlardan
yapılan yorum ve kullanımların dayanaklarıyla ilgilidir. Bir ölçüm tutarlı
biçimde yanlış yapılıyorsa yüksek güvenirlik gösterebilir; ancak amaçlanan
yapıyı doğru temsil etmeyebilir.

\begin{table}[htbp]
\centering
\caption{Güvenirlik ve geçerlik sorularının ayrımı}
\label{tab:reliability-validity}
\begin{tabular}{@{}p{0.23\textwidth}p{0.32\textwidth}p{0.32\textwidth}@{}}
\toprule
 & Güvenirlik & Geçerlik \\
\midrule
Temel soru & Puan ne kadar tutarlı ve hassastır? & Puanın amaçlanan yorumu ne ölçüde desteklenir? \\
Kanıt örneği & İç tutarlılık, zaman içi kararlılık, değerlendirici uyumu & İçerik, yapı, ölçüt ve sonuçlara ilişkin kanıtlar \\
Yanlış çıkarım & ``Alfa yüksekse ölçek geçerlidir.'' & ``Geçerli bir ölçek her grupta aynı güvenirliği verir.'' \\
Raporlama & Katsayı, örneklem, puanlama ve belirsizlik & Amaç, hedef grup ve çoklu kanıt kaynakları \\
\bottomrule
\end{tabular}
\end{table}

\begin{mistakebox}
Yüksek alfa tek boyutluluk kanıtı değildir. Birden fazla ilişkili alt boyut,
çok benzer maddeler veya yalnız madde sayısının artması da alfayı yükseltebilir
\citep{sijtsma2009,mcneish2018}.
\end{mistakebox}

\section{Hangi güvenirlik sorusu soruluyor?}

``Güvenirlik'' tek bir yöntem değildir. Katsayı, araştırmanın hata kaynağına
göre seçilmelidir.

\begin{figure}[htbp]
\centering
\resizebox{0.98\linewidth}{!}{%
\begin{tikzpicture}[>=stealth,
  box/.style={draw=PrismBlue!70,fill=PrismLight,rounded corners,align=center,
    minimum width=3.8cm,minimum height=1.05cm,font=\small},
  note/.style={font=\scriptsize,align=center,text=PrismGray}]
\node[box] (q) at (0,2.2) {Başlıca hata kaynağı nedir?};
\node[box] (items) at (-5,0) {Madde örneklemi\\İç tutarlılık};
\node[box] (time) at (0,0) {Zaman ve uygulama\\Test--tekrar test};
\node[box] (rater) at (5,0) {Değerlendirici\\Kappa veya ICC};
\draw[->,thick,PrismBlue] (q) -- (items);
\draw[->,thick,PrismBlue] (q) -- (time);
\draw[->,thick,PrismBlue] (q) -- (rater);
\node[note] at (-5,-1.25) {alfa, omega,\\madde analizi};
\node[note] at (0,-1.25) {kararlılık ve\\sistematik değişim};
\node[note] at (5,-1.25) {kategorik veya\\sürekli puanlama};
\end{tikzpicture}%
}
\caption{Araştırma sorusuna göre temel güvenirlik kanıtı seçimi.}
\label{fig:reliability-selection}
\end{figure}

\subsection{İç tutarlılık}

Aynı uygulamada, aynı puanı oluşturmak üzere birleştirilen maddelerin birlikte
değişimini inceler. Alfa ve omega bu başlık altında değerlendirilir. İç
tutarlılık, zaman içindeki kararlılığı veya değerlendiriciler arası uyumu
göstermez.

\subsection{Test--tekrar test güvenirliği}

Aynı ölçümün uygun bir zaman aralığıyla yeniden uygulanmasında puanların
kararlılığını inceler. Basit Pearson korelasyonu bireylerin sırasını özetler;
bütün puanların ikinci uygulamada sistematik olarak yükselmesini tek başına
yakalamayabilir. Amaç mutlak uyumsa uygun sınıf içi korelasyon katsayısı (ICC),
ortalama değişim ve fark grafikleri birlikte düşünülmelidir.

\subsection{Değerlendiriciler arası güvenirlik}

Kategorik kararlar için kappa ailesi, sürekli puanlar için uygun ICC modeli
kullanılabilir. Katsayı seçimi; değerlendiricilerin sabit mi rastgele mi
olduğu, tek değerlendirici mi ortalama puan mı kullanılacağı ve mutlak uyum mu
tutarlılık mı hedeflendiğine bağlıdır.

\section{Analizden önce veri ve puanlama denetimi}

Güvenirlik katsayısı hesaplanmadan önce şu kontroller yapılmalıdır:

\begin{enumerate}
  \item Her madde için geçerli puan aralığı ve eksik değer kodları doğrulanır.
  \item Olumsuz maddeler puanlama anahtarına göre ters çevrilir.
  \item Toplam veya ortalama puanın kaç geçerli maddeyle hesaplanacağı belirlenir.
  \item Maddelerin aynı hedef puanda birleştirilmesinin kuramsal gerekçesi incelenir.
  \item Madde dağılımları, tavan--taban yığılmaları ve olağandışı örüntüler görülür.
  \item Alt boyutlu bir yapı varsa bütün maddeler için tek katsayı vermek yerine
  alt boyutlar ayrı değerlendirilir.
\end{enumerate}


Beşli bir maddede ters puanlama
\[
X_{\text{ters}}=6-X
\]
ile yapılır. Genel olarak en küçük değer $L$, en büyük değer $U$ ise
\[
X_{\text{ters}}=L+U-X
\]
kullanılır.

\begin{mistakebox}
Madde metni olumsuz diye otomatik ters puanlama yapılmaz. Ölçeğin özgün
puanlama anahtarı incelenmelidir. Bazı olumsuz ifadeler zaten yüksek yapı
düzeyini gösterebilir veya farklı bir alt boyuta ait olabilir.
\end{mistakebox}

\section{Cronbach alfa katsayısı}
\index{Cronbach alfa}

$k$ maddelik bir toplam puan için alfa
\[
\alpha=\frac{k}{k-1}
\left(1-\frac{\sum_{j=1}^{k}s_j^2}{s_X^2}\right)
\]
biçiminde hesaplanır. Burada $s_j^2$ her maddenin varyansı, $s_X^2$ ise madde
toplamının varyansıdır. Standartlaştırılmış maddelerde ortalama maddeler arası
korelasyon $\bar r$ kullanılarak
\[
\alpha_{\text{std}}=\frac{k\bar r}{1+(k-1)\bar r}
\]
elde edilir \citep{cronbach1951}.

Formül iki önemli özelliği görünür kılar: Maddeler birlikte değiştikçe alfa
yükselir; aynı ortalama ilişki düzeyinde madde sayısı arttıkça da alfa
yükselme eğilimindedir. Bu nedenle farklı uzunluktaki ölçeklerin alfalarını
yalnız sayısal büyüklükle karşılaştırmak yanıltıcı olabilir.

\begin{assumptionbox}
Alfanın güvenirliğe yakın yorumlanması; maddelerin aynı yapıyı uygun biçimde
temsil etmesi, hata yapısının sorunlu olmaması ve maddelerin hedef puana
katkılarının aşırı farklılaşmaması gibi koşullara dayanır. Alfa bu koşulları
kendi başına sınamaz.
\end{assumptionbox}

\section{Adım adım ters madde örneği}

Tablo~\ref{tab:reliability-small-data}, yalnız hesaplama mantığını göstermek
için hazırlanmış 12 kişilik küçük bir öğretim verisidir. Gerçek bir ölçek
çalışması için bu örneklem yeterli değildir.

\begin{table}[htbp]
\centering
\caption{Beş maddelik öğretim verisi; M4 ters yönde yazılmıştır}
\label{tab:reliability-small-data}
\small
\begin{tabular}{@{}crrrrr@{}}
\toprule
Kişi & M1 & M2 & M3 & M4 & M5 \\
\midrule
1 & 1 & 2 & 1 & 5 & 2 \\
2 & 2 & 2 & 2 & 4 & 2 \\
3 & 2 & 3 & 2 & 4 & 3 \\
4 & 3 & 2 & 3 & 3 & 3 \\
5 & 3 & 3 & 4 & 2 & 3 \\
6 & 4 & 3 & 4 & 2 & 4 \\
7 & 4 & 4 & 3 & 2 & 4 \\
8 & 5 & 4 & 5 & 1 & 4 \\
9 & 4 & 5 & 4 & 1 & 5 \\
10 & 5 & 5 & 5 & 1 & 5 \\
11 & 3 & 4 & 3 & 2 & 4 \\
12 & 2 & 3 & 3 & 3 & 2 \\
\bottomrule
\end{tabular}
\end{table}

M4 değiştirilmeden hesaplanan alfa $0{,}294$'tür. Madde
$M4_{\text{ters}}=6-M4$ biçiminde düzeltildiğinde alfa $0{,}964$ olur. Büyük
değişim, ilk sonucun ölçeğin niteliğinden çok puanlama hatasını yansıttığını
gösterir.

\begin{table}[htbp]
\centering
\caption{Ters puanlama sonrasında madde tanıları}
\label{tab:reliability-item-diagnostics}
\begin{tabular}{@{}lrrr@{}}
\toprule
Madde & Ortalama & Düzeltilmiş madde--toplam $r$ & Madde silinirse $\alpha$ \\
\midrule
M1 & 3,17 & 0,934 & 0,949 \\
M2 & 3,33 & 0,831 & 0,966 \\
M3 & 3,25 & 0,880 & 0,958 \\
M4 ters & 3,50 & 0,965 & 0,944 \\
M5 & 3,42 & 0,890 & 0,957 \\
\bottomrule
\end{tabular}
\end{table}

Örnekte bütün düzeltilmiş madde--toplam korelasyonları pozitiftir. M2
silindiğinde alfa çok küçük bir miktar yükselmektedir. Bu fark M2'yi otomatik
olarak silmek için yeterli değildir. Maddenin kuramsal kapsamı, anlaşılabilirliği,
benzer maddelerle fazlalığı ve başka örneklemlerdeki davranışı incelenmelidir.
Ayrıca $\alpha\approx0{,}96$ maddelerin aşırı benzer olabileceği olasılığını da
gündeme getirir.

\section{Madde--toplam korelasyonu nasıl okunur?}
\index{madde--toplam korelasyonu}

Düzeltilmiş madde--toplam korelasyonu, bir madde ile o madde çıkarılarak
hesaplanan kalan toplam arasındaki ilişkiyi özetler. Maddenin kendi toplamına
dahil edildiği düzeltilmemiş korelasyon yapay olarak yüksek olabilir.

Negatif bir değer şu olasılıkları düşündürür:

\begin{itemize}
  \item ters maddenin yanlış kodlanması,
  \item veri giriş veya eksik değer kodlama hatası,
  \item maddenin farklı bir yapıyı ölçmesi,
  \item belirsiz veya çift anlamlı madde ifadesi,
  \item çok küçük ve kararsız örneklem.
\end{itemize}

Tek bir eşik bütün ölçekler için karar kuralı değildir. Madde istatistiği,
içerik geçerliği ve ölçeğin hedeflediği kapsamla birlikte değerlendirilmelidir.

\section{Omega neden raporlanabilir?}
\index{omega katsayısı}

Alfa maddelerin ortak yapıya katkılarının benzer olmasına güçlü biçimde
dayanabilir. Tek faktörlü eşlenik bir modelde yükler $\lambda_j$ ve özgül hata
varyansları $\theta_j$ ise toplam omega basitleştirilmiş olarak
\[
\omega=\frac{\left(\sum_{j=1}^{k}\lambda_j\right)^2}
{\left(\sum_{j=1}^{k}\lambda_j\right)^2+\sum_{j=1}^{k}\theta_j}
\]
biçiminde ifade edilir. Omega, maddelerin faktör yüklerinin eşit olmasını
gerektirmeyen model temelli bir yaklaşım sunar \citep{dunn2014,mcneish2018}.

\begin{sezgibox}
Alfa gözlenen madde varyans ve kovaryanslarından tek bir özet üretir. Omega
ise maddelerin ortak yapıya farklı güçlerde bağlanabileceğini açıkça modele
katar. Bununla birlikte yanlış faktör modeliyle hesaplanan omega da güvenilir
bir sonuç garanti etmez.
\end{sezgibox}

Alfa ile omega birbirine yakınsa bu, her varsayımın doğrulandığı anlamına
gelmez; çok farklıysa puan yapısı, madde yükleri ve hata ilişkileri daha
dikkatle incelenmelidir. Alt boyutlu ölçeklerde toplam omega yanında alt ölçek
puanlarının anlamlılığı da değerlendirilir.

\section{Sık kullanılan eşikler neden yeterli değildir?}

``0,70'in üzeri kabul edilir'' gibi sabit kurallar araştırma kararının yerini
tutmaz. Katsayının yeterliliği en az şu unsurlara bağlıdır:

\begin{itemize}
  \item kararın bireysel mi grup düzeyinde mi olduğu,
  \item yanlış kararın uygulamadaki sonucu,
  \item ölçeğin uzunluğu ve kapsamı,
  \item örneklemin heterojenliği,
  \item katsayının belirsizliği,
  \item başka güvenirlik ve geçerlik kanıtları.
\end{itemize}

Araştırma gruplarını betimlemek için kullanılan bir puan ile bireysel klinik
karara temel olacak puanın gerektirdiği hassasiyet aynı değildir. Katsayı
yorumunda kullanım amacı mutlaka belirtilmelidir.

\section{Eksik veri, uç yanıt ve örneklem etkisi}

Liste bazında silme, yalnız bütün maddeleri yanıtlayan kişileri analize alır ve
yanıtlamama sistematikse sonucu yanlılaştırabilir. Çift bazında kovaryans
hesapları ise farklı hücrelerin farklı katılımcılara dayanmasına yol açabilir.
Kullanılan eksik veri kuralı ve analize giren kişi sayısı raporlanmalıdır.

Örneklem çok homojense gerçek puan değişkenliği azalabilir ve güvenirlik
katsayısı düşebilir. Aşırı heterojen örneklem katsayıyı yükseltebilir. Bu
nedenle başka bir çalışmada bulunan alfa, yeni örneklem için otomatik olarak
aktarılmaz.

\section{Python ile alfa ve madde tanıları}

\begin{lstlisting}[style=prismPython,caption={Cronbach alfa ve madde tanıları için Python kodu}]
import numpy as np
import pandas as pd

def cronbach_alpha(frame):
    x = frame.to_numpy(dtype=float)
    k = x.shape[1]
    item_variances = x.var(axis=0, ddof=1).sum()
    total_variance = x.sum(axis=1).var(ddof=1)
    return k / (k - 1) * (1 - item_variances / total_variance)

items = ["M1", "M2", "M3", "M4", "M5"]
df = pd.read_csv("guvenirlik.csv")

# Besli maddede ters puanlama
df["M4_t"] = 6 - df["M4"]
scored = df[["M1", "M2", "M3", "M4_t", "M5"]]

print("alpha =", cronbach_alpha(scored))

for item in scored.columns:
    remaining = scored.drop(columns=item)
    corrected_total = remaining.sum(axis=1)
    item_total = scored[item].corr(corrected_total)
    alpha_deleted = cronbach_alpha(remaining)
    print(item, item_total, alpha_deleted)
\end{lstlisting}

Kod küçük öğretim verisinde $\alpha=0{,}963643$ üretir. Gerçek analizde
eksik değer işlemi, madde aralığı ve ters puanlama koddan önce belgelenmelidir.

\section{R ile alfa ve omega}

Temel alfa hesabı ek paket olmadan yapılabilir:

\begin{lstlisting}[style=prismR,caption={R ile alfa hesabı ve omega için model temelli seçenek}]
cronbach_alpha <- function(x) {
  x <- as.data.frame(x)
  k <- ncol(x)
  k / (k - 1) *
    (1 - sum(sapply(x, var, na.rm = TRUE)) /
       var(rowSums(x), na.rm = TRUE))
}

d <- read.csv("guvenirlik.csv")
d$M4_t <- 6 - d$M4
puanlanan <- d[c("M1", "M2", "M3", "M4_t", "M5")]
cronbach_alpha(puanlanan)

# psych paketi kuruluysa ayrintili alfa ve omega:
# psych::alpha(puanlanan, check.keys = FALSE)
# psych::omega(puanlanan, nfactors = 1)
\end{lstlisting}

Otomatik anahtar düzeltme seçenekleri öğretim amacıyla yararlı görünse de,
ters maddelerin yazılım tarafından tahmin edilmesi özgün puanlama anahtarının
yerini tutmaz.

\section{SPSS ile uygulama ve çıktı okuma}

\begin{lstlisting}[style=prismSPSS,caption={SPSS ile ters puanlama ve alfa analizi}]
RECODE M4 (1=5) (2=4) (3=3) (4=2) (5=1) INTO M4_T.
EXECUTE.

RELIABILITY
  /VARIABLES=M1 M2 M3 M4_T M5
  /SCALE('Yardim arama tutumu') ALL
  /MODEL=ALPHA
  /STATISTICS=DESCRIPTIVE SCALE CORR.
\end{lstlisting}

\begin{outputguidebox}
SPSS çıktısında en az şu bölümler okunmalıdır: (1) analize giren ve dışlanan
kişi sayısı, (2) ölçek için alfa, (3) madde ortalama ve standart sapmaları,
(4) düzeltilmiş madde--toplam korelasyonları, (5) madde silinirse alfa.
Yalnız ``Reliability Statistics'' tablosundaki tek katsayıyı raporlamak
yeterli değildir.
\end{outputguidebox}

SPSS'in standart \texttt{RELIABILITY} komutu alfa verir; omega için faktör
modeli kurabilen ek yazılım veya uygun bir R/Python uygulaması gerekir.

\section{Bilimsel raporlama örneği}

Güvenirlik paragrafı; örneklemi, madde sayısını, puanlama işlemini, katsayıyı,
madde tanılarını ve yorum sınırını birlikte vermelidir:

\begin{quote}
``Beş maddelik yardım arama tutumu puanı 120 üniversite öğrencisinde
incelenmiştir. Olumsuz yöndeki dördüncü madde özgün anahtara göre ters
puanlanmıştır. İç tutarlılık için Cronbach alfa $0{,}84$ ve tek faktörlü
modelden toplam omega $0{,}86$ bulunmuştur. Düzeltilmiş madde--toplam
korelasyonları $0{,}46$ ile $0{,}71$ arasındadır. Sonuçlar bu örneklemdeki
toplam puanın iç tutarlılığını desteklemektedir; tek boyutluluk, zaman içi
kararlılık veya geçerlik kanıtı olarak yorumlanmamıştır.''
\end{quote}

Katsayı için güven aralığı elde edilebiliyorsa nokta tahminiyle birlikte
verilmelidir. Örneklem küçükse katsayının kararsızlığı özellikle belirtilir.

\section{Bir güvenirlik araştırmasını değerlendirirken}

\begin{enumerate}
  \item Puanın hangi amaç ve hedef grup için kullanılacağı açık mı?
  \item Maddelerin aynı toplamda birleştirilmesi kuramsal ve ampirik olarak savunulmuş mu?
  \item Ters maddeler ve eksik değerler doğru işlenmiş mi?
  \item Katsayı araştırma sorusundaki hata kaynağına uygun mu?
  \item Alfa tek boyutluluk veya geçerlik kanıtı gibi sunulmuş mu?
  \item Madde silme kararları yalnız katsayı artışına mı dayanıyor?
  \item Katsayının örnekleme bağlı olduğu ve belirsizliği tartışılmış mı?
  \item İç tutarlılığın yanında gerekliyse zaman veya değerlendirici kanıtı verilmiş mi?
\end{enumerate}

\section{Bölüm sonu çözümlü problemler}

\begin{enumerate}
  \item \textbf{Standartlaştırılmış alfa.} Altı maddelik bir ölçekte ortalama
  maddeler arası korelasyon $\bar r=0{,}30$ bulunmuştur. Standartlaştırılmış
  alfayı hesaplayın.

  \textbf{Çözüm:}
  \[
  \alpha_{\mathrm{std}}
  =\frac{6(0{,}30)}{1+5(0{,}30)}
  =\frac{1{,}80}{2{,}50}=0{,}72.
  \]
  Bu değer tek başına ölçeğin tek boyutlu veya geçerli olduğunu göstermez.

  \item \textbf{Ters madde hatası.} Beş maddelik bir ölçeğin alfası 0,31'dir.
  Dördüncü maddenin diğer maddelerle korelasyonları negatiftir ve madde
  olumsuz yönde yazılmıştır. Araştırmacının izlemesi gereken sırayı yazın.

  \textbf{Çözüm:} Önce özgün puanlama anahtarı incelenir; yalnız madde metnine
  bakılarak karar verilmez. Madde gerçekten tersse doğru aralık formülüyle
  yeniden kodlanır. Değer aralığı ve frekanslar kontrol edilir; ardından alfa
  ve düzeltilmiş madde--toplam korelasyonları yeniden hesaplanır. İlk katsayı
  ölçeğin zayıflığı olarak değil, olası puanlama hatasının işareti olarak ele
  alınır.

  \item \textbf{Madde silme kararı.} Toplam alfa 0,82, bir madde silindiğinde
  alfa 0,83'tür. Maddenin düzeltilmiş madde--toplam korelasyonu 0,38'dir.
  Madde silinmeli midir?

  \textbf{Çözüm:} Yalnız 0,01'lik alfa artışı silme kararı için yeterli
  değildir. Maddenin ölçek kapsamındaki özgün katkısı, ifadesinin açıklığı,
  alt boyut yapısı, başka örneklemlerdeki davranışı ve silmenin içerik
  geçerliğine etkisi incelenmelidir. Madde kuramsal olarak önemliyse korunması
  daha uygun olabilir.

  \item \textbf{Uygun kanıtı seçme.} Üç danışman aynı görüşme kayıtlarını
  bağımsız olarak 0--4 arasında puanlamaktadır. Araştırma ekibi yalnız
  Cronbach alfa raporlamıştır. Sorunu açıklayın.

  \textbf{Çözüm:} Temel hata kaynağı madde örneklemi değil değerlendiricidir.
  Sürekli/sıralı puanın kullanımına ve değerlendirici tasarımına uygun bir ICC
  modeli veya uygun uyum katsayısı seçilmelidir. Model türü, tek değerlendirici
  ya da ortalama puan ve mutlak uyum/tutarlılık hedefi açıkça belirtilmelidir.
\end{enumerate}

\bolumkaynaklari{\cite{cronbach1951,sijtsma2009,dunn2014,mcneish2018,appelbaum2018}}

\section*{Hatırla--Uygula--Yansıt}

\begin{enumerate}
\renewcommand{\labelenumi}{\textbf{\Alph{enumi}.}}
  \item \textbf{Hatırla:} Güvenirliği neden ölçeğin değişmez özelliği olarak tanımlamamak gerekir?
  \item \textbf{Ayır:} İç tutarlılık, test--tekrar test ve değerlendiriciler arası güvenirliğe birer araştırma sorusu yazın.
  \item \textbf{Kodla:} 0--4 aralığındaki ters bir madde için dönüşüm formülünü yazın.
  \item \textbf{Hesapla:} $k=6$ ve ortalama maddeler arası korelasyon $0{,}30$ ise standartlaştırılmış alfayı bulun.
  \item \textbf{Eleştir:} ``Alfa 0,91 olduğu için ölçek tek boyutlu ve geçerlidir'' cümlesini düzeltin.
  \item \textbf{Tanıla:} Bir maddenin düzeltilmiş madde--toplam korelasyonu negatifse hangi kontrolleri yaparsınız?
  \item \textbf{Karar ver:} Bir madde silindiğinde alfa 0,80'den 0,81'e çıkıyorsa hangi ek kanıtları incelersiniz?
  \item \textbf{Karşılaştır:} Alfa ile omega arasındaki temel model farkını açıklayın.
  \item \textbf{Uygula:} Kendi veri setinizde madde dağılımlarını, alfa ve madde tanılarını Python, R veya SPSS ile üretin.
  \item \textbf{Raporla:} Örneklem, madde sayısı, puanlama, katsayı, madde aralığı ve yorum sınırı içeren kısa paragraf yazın.
\end{enumerate}

\bolumbaglantisi{chapters/pdr/16-olcek-guvenirligi.tex}